\section{ORDER OF BATTLE}

The starting forces, their location and their strength (in parenthesis) are listed below:

\subsection{Union}

The Union forces are divided and caught off guard by Jackson’s rapid advance.
Strengths are adjusted to reflect the state of the campaign as of May 23, 1862.

\subsubsection{Department of the Shenandoah}

\begin{tabular}{ |l|c|c|c| }
  \hline
  Unit & Strength & Location & Type \\
  \hline
  Banks & 1 & Strasburg & HQ \\
  Taylor & 2 & Strasburg & Inf \\
  Donnelly & 2 & Strasburg & Inf \\
  Hatch & 2 & Strasburg & Cav \\
  Gordon & 2 & Strasburg & Inf \\
  Kenly & 1 & Front Royal & Inf \\
  \hline
\end{tabular}

\subsubsection{Mountain Department}

\begin{tabular}{ |l|c|c|c| }
  \hline
  Unit & Strength & Location & Type \\
  \hline
  Fremont & 1 & Franklin & HQ \\
  DeBeck & 3 & Franklin & Art \\
  Cluseret & 2 & Franklin & Inf \\
  Milroy & 2 & Cave & Inf \\
  Schenk & 2 & Cave & Inf \\
  Bohlen & 2 & Moorefield & Inf \\
  Koltes & 2 & Moorefield & Inf \\
  Stahel & 3 & Purgitsville & Inf \\
  \hline
\end{tabular}

\textit{\textbf{Designer’s Note:} The precise locations of Fremont’s forces are not strictly accurate. Rather, they are a game concession meant to reflect Fremont’s somewhat ponderous response when ordered to march on Harrisonburg. For this scenario Fremont will have to spend a little time organizing his forces and then decide where to press. Historically he decided the roads were too difficult, and marched north to Moorefield and then east.}
